% !TEX root = ../main.tex

%----------------------------------------------------------------------
\subsection{Defect engineering in two-dimensional materials}
%----------------------------------------------------------------------

Point defects---substitutional impurities, interstitials, and
adsorbates---fundamentally govern the electronic, optical, and catalytic
properties of two-dimensional (2D) materials~\cite{imp2d2023}.
In transition-metal dichalcogenides (TMDs) such as MoS$_2$ and WSe$_2$,
a single dopant atom can shift the Fermi level, create mid-gap trap states,
or activate catalytic edge sites---properties directly relevant to
photocatalytic water splitting, single-photon emission, and
next-generation field-effect transistors.
The defect formation energy $\ef$ quantifies the thermodynamic cost of
introducing such a defect and serves as the primary descriptor for
high-throughput defect screening: low-$\ef$ dopants are thermodynamically
accessible under typical growth conditions, while high-$\ef$ species
indicate kinetically trapped or unstable configurations.

However, the combinatorial space of possible host--dopant configurations
is enormous.
The Impurities in 2D Materials Database (IMP2D)~\cite{imp2d2023} covers 44
monolayer hosts doped with 65 elements, yielding a $44\times65 = 2{,}860$
matrix of which only 91.2\% entries have converged DFT results
(10\,641 configurations).
Each DFT supercell calculation requires hours of CPU time, making exhaustive
enumeration impractical for the remaining gaps and for extensions to larger
supercells or additional host families.

%----------------------------------------------------------------------
\subsection{Limitations of existing approaches}
%----------------------------------------------------------------------

Machine-learning (ML) surrogate models offer orders-of-magnitude speedup
over DFT.
The Atomistic Line Graph Neural Network (ALIGNN)~\cite{alignn2021}, a
state-of-the-art crystal property predictor with 4.03\,M parameters, achieves
an MAE of 0.540\eV{} on IMP2D under random 80/10/10 splitting.
El~Alouani \textit{et al.}~\cite{elalouani2024} applied gradient-boosted
trees with hand-crafted Jacobi--Legendre descriptors, reaching 0.518\eV{}.
In the 3D bulk domain, defect-specific GNNs have emerged:
DefiNet~\cite{definet2025} introduces equivariant GNN with explicit defect
markers for 14\,866 bulk defects;
Fang and Yan~\cite{fang2025} leverage persistent homology features to
reduce defect $\ef$ MAE by 55\%;
and Kiyohara \textit{et al.}~\cite{kiyohara2025} predict charged-defect
energies from crystal structure alone.
However, none of these approaches address the 2D setting or provide
systematic out-of-distribution (OOD) evaluation with prospective validation.

Despite these advances, four critical gaps persist:
\begin{enumerate}
  \item \textbf{Parameter efficiency.}
    ALIGNN's 4.03\,M parameters are disproportionate for a dataset of
    ${\sim}10$\,k samples, risking overfitting.
    Our scaling analysis reveals $\alpha_\mathrm{data} = -0.40$ vs.\
    $\beta_\mathrm{params} \approx -0.01$ ($R^2 = 0.95$)---data, not
    model capacity, is the bottleneck.

  \item \textbf{Physical inductive bias for defects.}
    Standard GNNs are bounded by a fixed cutoff radius, yet point defects in
    2D materials induce strain fields and charge-density perturbations
    extending well beyond typical $r_\mathrm{cut}$~\cite{fang2025}.
    Moreover, the readout stage typically uses mean pooling, which dilutes the
    extreme features of the defect atom among dozens of host atoms.

  \item \textbf{Rigorous innovation assessment.}
    Most architecture papers report only point estimates of accuracy gains.
    With typical materials test sets of ${\sim}10^3$ samples, statistical
    power is limited; bootstrap hypothesis testing is essential to
    distinguish genuine improvements from noise.

  \item \textbf{Prospective validation.}
    All prior work on IMP2D reports only random-split accuracy---a measure
    of interpolation.
    No study has evaluated whether predictions transfer to \emph{unseen
    host materials} or validated model-recommended candidates with
    independent DFT calculations.
\end{enumerate}

%----------------------------------------------------------------------
\subsection{Our contributions}
%----------------------------------------------------------------------

We address these gaps by introducing \dart{} (Defect-Aware Radial
Transformer), a compact model with three architectural innovations and
twelve systematically evaluated physical inductive biases, accompanied by
rigorous statistical testing and multi-tier generalization assessment.

Our central thesis is that, in the data-limited regime of materials science
($N \sim 10^4$), \emph{physically informed architecture design combined with
careful training engineering} yields larger gains than scaling model capacity
or adding auxiliary loss terms.
\dart{} demonstrates this: none of six attempted auxiliary losses (label
reweighting, contrastive, focal, heteroscedastic) improve over the base
model, while targeted architectural changes (gated pooling, environment
enrichment) provide consistent gains.

Specifically, our contributions are:

\begin{enumerate}
  \item \textbf{\dart{} architecture} (\S\ref{sec:arch}):
    A 0.83\,M-parameter model combining local bond--angle message passing
    with global periodic self-attention, augmented by three innovations:
    \emph{defect-aware gated attention pooling} ($-19.9\%$),
    \emph{local environment enrichment} ($-18.0\%$), and
    \emph{pre-norm residual connections} ($-21.0\%$).
    Together, these reduce single-model MAE from 0.516 to 0.381\eV{}
    ($-26.2\%$), with the best variant (MoE readout) reaching
    \textbf{0.379\eV{}} ($-30\%$ vs.\ ALIGNN).

  \item \textbf{Systematic innovation ablation} (\S\ref{sec:ablation_innov}):
    Twelve physically motivated innovations---spanning architecture,
    loss functions, and training strategies---evaluated with bootstrap
    hypothesis testing ($n = 10{,}000$ resamples, 95\% CI).
    Six match the baseline ($p > 0.05$); six are significantly harmful
    ($p < 0.05$), providing actionable negative results.

  \item \textbf{Efficient ensemble} (\S\ref{sec:ensemble}):
    Greedy forward selection from 20 models identifies an 8-model
    ensemble at \textbf{0.344\eV{}} ($-36\%$ vs.\ ALIGNN), surpassing
    a prior 29-model ensemble (0.359\eV{}) with 3.6$\times$ fewer members.

  \item \textbf{Graduated OOD assessment and calibrated UQ}
    (\S\ref{sec:ood}, \S\ref{sec:uq}):
    Three-tier protocol (ID $\to$ intra-family LOHO $\to$ compositional
    block-out) quantifying deployment boundaries, with temperature-scaled
    ensemble uncertainty achieving 93.4\% coverage at 90\% nominal level.

  \item \textbf{Prospective DFT validation} (\S\ref{sec:dft}):
    Independent first-principles verification of model-recommended
    candidates, bridging the gap between benchmark accuracy and
    real-world deployment.
    % TODO: expand when DFT results are finalized
\end{enumerate}
